% GENERATED FILE. Edit canonical lessons, then run npm run generate:books.
% canonical-source-hash: fnv1a64:90d79ef16e1e1394
% canonical-lessons: TA-C390-cirtiruttam, TA-C390-munnerram, TA-C390-mempatuttu, TA-C390-ciramai, TA-C390-valuppatuttu

\chapter{Change for the Better}
\label{ch:ta-change-for-the-better}
\hlchaptermodality{390}

\begin{chapteropening}
\textbf{By the end of this chapter:} \emph{I can talk about reform, progress and improvement.}

Complete the last lesson of chapter 390: I can talk about reform, progress and improvement.
\end{chapteropening}

\section[\texorpdfstring{\ta{சீர்திருத்தம்} (cīrtiruttam)}{சீர்திருத்தம் (cīrtiruttam)}]{\ta{சீர்திருத்தம்} (cīrtiruttam) — a reform}

\label{lesson:TA-C390-cirtiruttam}



\begin{quote}
Before the new one: say the Tamil for a course of action, then the Tamil for to administer.
\end{quote}

\subsection*{You'll want to know: \ta{சீர்திருத்தம்}}
\textbf{\ta{சீர்திருத்தம்}} — \emph{cīrtiruttam} — \textquotedblleft{}a reform\textquotedblright{}.

\begin{quote}
\ta{கல்வியில்} \ta{சீர்திருத்தம்} \ta{வந்தால்} \ta{எல்லோருக்கும்} \ta{நல்லது}.
\end{quote}

\begin{quote}
\emph{kalviyil cīrtiruttam vantāl ellōrukkum nallatu.}
\end{quote}

\begin{quote}
If education is reformed, it will be good for everyone.
\end{quote}

Plural: \ta{சீர்திருத்தங்கள்}, \emph{cīrtiruttaṅkaḷ}.

\ta{சீர்}, good order + \ta{திருத்தம்}, correction. Plural \ta{சீர்திருத்தங்கள்}.

\begin{cousinweb}
\textbf{\ta{சீர்திருத்தம்}} (\emph{cīrtiruttam}) is native: \emph{cīr}, \textquotedblleft{}to be excellent, fitting\textquotedblright{} (DEDR 2624), and \emph{tiruttam}, \textquotedblleft{}correction\textquotedblright{}, from \emph{tiruttu}, \textquotedblleft{}to correct, reform\textquotedblright{} (DEDR 3251).
\end{cousinweb}

\begin{grammarlens}[title={What you've built}]
One more word for saying what happens if.
\end{grammarlens}

\subsection*{Your turn}
\begin{itemize}
\raggedright
  \item \emph{Say it:} \emph{cīrtiruttam}
  \item \emph{Say it:} \emph{cīrtiruttam}, once more
  \item \emph{Say it:} \emph{kalviyil cīrtiruttam vantāl ellōrukkum nallatu.} — If education is reformed, it will be good for everyone.
  \item \emph{Recall:} say the Tamil for a course of action, then the Tamil for to administer, then say \emph{cīrtiruttam} again
\end{itemize}

\begin{tcolorbox}[breakable,colback=teal!4,colframe=teal!35!black,title={Before you move on}]
What does \ta{சீர்திருத்தம்} mean? (\textquotedblleft{}A reform\textquotedblright{}.) Say it once more.
\end{tcolorbox}
\section[\texorpdfstring{\ta{முன்னேற்றம்} (muṉṉēṟṟam)}{முன்னேற்றம் (muṉṉēṟṟam)}]{\ta{முன்னேற்றம்} (muṉṉēṟṟam) — progress}

\label{lesson:TA-C390-munnerram}



\begin{quote}
Before the new one: say the Tamil for to administer, then the Tamil for a reform.
\end{quote}

\subsection*{You'll want to know: \ta{முன்னேற்றம்}}
\textbf{\ta{முன்னேற்றம்}} — \emph{muṉṉēṟṟam} — \textquotedblleft{}progress\textquotedblright{}.

\begin{quote}
\ta{தினமும்} \ta{பயிற்சி} \ta{செய்தால்} \ta{நல்ல} \ta{முன்னேற்றம்} \ta{கிடைக்கும்}.
\end{quote}

\begin{quote}
\emph{tiṉamum payiṟci ceytāl nalla muṉṉēṟṟam kiṭaikkum.}
\end{quote}

\begin{quote}
If you practise every day, you will make good progress.
\end{quote}

\ta{முன்னேற்றம்} \ta{அடை}, to make progress.

\begin{cousinweb}
\textbf{\ta{முன்னேற்றம்}} (\emph{muṉṉēṟṟam}) is native: \emph{muṉ}, \textquotedblleft{}front, before\textquotedblright{} (DEDR 5020), and \emph{ēṟṟam}, \textquotedblleft{}mounting, rising\textquotedblright{}, from \emph{ēṟu}, \textquotedblleft{}to rise\textquotedblright{} (DEDR 916): a rising forward.
\end{cousinweb}

\begin{grammarlens}[title={What you've built}]
One more word for saying what happens if.
\end{grammarlens}

\subsection*{Your turn}
\begin{itemize}
\raggedright
  \item \emph{Say it:} \emph{muṉṉēṟṟam}
  \item \emph{Say it:} \emph{muṉṉēṟṟam}, once more
  \item \emph{Say it:} \emph{tiṉamum payiṟci ceytāl nalla muṉṉēṟṟam kiṭaikkum.} — If you practise every day, you will make good progress.
  \item \emph{Recall:} say the Tamil for to administer, then the Tamil for a reform, then say \emph{muṉṉēṟṟam} again
\end{itemize}

\begin{tcolorbox}[breakable,colback=teal!4,colframe=teal!35!black,title={Before you move on}]
What does \ta{முன்னேற்றம்} mean? (\textquotedblleft{}Progress\textquotedblright{}.) Say it once more.
\end{tcolorbox}
\section[\texorpdfstring{\ta{மேம்படுத்து} (mēmpaṭuttu)}{மேம்படுத்து (mēmpaṭuttu)}]{\ta{மேம்படுத்து} (mēmpaṭuttu) — to improve; to upgrade}

\label{lesson:TA-C390-mempatuttu}



\begin{quote}
Before the new one: say the Tamil for a reform, then the Tamil for progress.
\end{quote}

\subsection*{You'll want to know: \ta{மேம்படுத்து}}
\textbf{\ta{மேம்படுத்து}} — \emph{mēmpaṭuttu} — \textquotedblleft{}to improve; to upgrade\textquotedblright{}.

\begin{quote}
\ta{சாலைகளை} \ta{மேம்படுத்தினால்} \ta{விபத்துகள்} \ta{குறையும்}.
\end{quote}

\begin{quote}
\emph{cālaikaḷai mēmpaṭuttiṉāl vipattukaḷ kuṟaiyum.}
\end{quote}

\begin{quote}
If the roads are improved, there will be fewer accidents.
\end{quote}

Past (I …): \ta{மேம்படுத்தினேன்}, \emph{mēmpaṭuttiṉēṉ}.

What is improved takes -\ta{ஐ}. The noun is \ta{மேம்பாடு}, improvement.

\begin{cousinweb}
\textbf{\ta{மேம்படுத்து}} (\emph{mēmpaṭuttu}) is native: \emph{mēm}, a form of \emph{mēl}, \textquotedblleft{}that which is above, superior\textquotedblright{} (DEDR 5086), with \emph{paṭuttu}, the causative of \emph{paṭu}: to raise higher.
\end{cousinweb}

\begin{grammarlens}[title={What you've built}]
One more word for saying what happens if.
\end{grammarlens}

\subsection*{Your turn}
\begin{itemize}
\raggedright
  \item \emph{Say it:} \emph{mēmpaṭuttu}
  \item \emph{Say it:} \emph{mēmpaṭuttu}, once more
  \item \emph{Say it:} \emph{cālaikaḷai mēmpaṭuttiṉāl vipattukaḷ kuṟaiyum.} — If the roads are improved, there will be fewer accidents.
  \item \emph{Recall:} say the Tamil for a reform, then the Tamil for progress, then say \emph{mēmpaṭuttu} again
\end{itemize}

\begin{tcolorbox}[breakable,colback=teal!4,colframe=teal!35!black,title={Before you move on}]
What does \ta{மேம்படுத்து} mean? (\textquotedblleft{}To improve; to upgrade\textquotedblright{}.) Say it once more.
\end{tcolorbox}
\section[\texorpdfstring{\ta{சீரமை} (cīramai)}{சீரமை (cīramai)}]{\ta{சீரமை} (cīramai) — to restructure; to put in order}

\label{lesson:TA-C390-ciramai}



\begin{quote}
Before the new one: say the Tamil for progress, then the Tamil for to improve.
\end{quote}

\subsection*{You'll want to know: \ta{சீரமை}}
\textbf{\ta{சீரமை}} — \emph{cīramai} — \textquotedblleft{}to restructure; to put in order\textquotedblright{}.

\begin{quote}
\ta{புதிய} \ta{மேலாளர்} \ta{அலுவலகத்தைச்} \ta{சீரமைத்தார்}.
\end{quote}

\begin{quote}
\emph{putiya mēlāḷar aluvalakattaic cīramaittār.}
\end{quote}

\begin{quote}
The new manager reorganised the office.
\end{quote}

Past (I …): \ta{சீரமைத்தேன்}, \emph{cīramaittēṉ}.

\ta{சீர்}, good order + \ta{அமை}, to arrange.

\begin{cousinweb}
\textbf{\ta{சீரமை}} (\emph{cīramai}) is native: \emph{cīr}, \textquotedblleft{}excellence, fitness\textquotedblright{} (DEDR 2624), and \emph{amai}, \textquotedblleft{}to be fitted, set in place\textquotedblright{} (DEDR 162): to set in good order.
\end{cousinweb}

\begin{grammarlens}[title={What you've built}]
One more word for saying what happens if.
\end{grammarlens}

\subsection*{Your turn}
\begin{itemize}
\raggedright
  \item \emph{Say it:} \emph{cīramai}
  \item \emph{Say it:} \emph{cīramai}, once more
  \item \emph{Say it:} \emph{putiya mēlāḷar aluvalakattaic cīramaittār.} — The new manager reorganised the office.
  \item \emph{Recall:} say the Tamil for progress, then the Tamil for to improve, then say \emph{cīramai} again
\end{itemize}

\begin{tcolorbox}[breakable,colback=teal!4,colframe=teal!35!black,title={Before you move on}]
What does \ta{சீரமை} mean? (\textquotedblleft{}To restructure; to put in order\textquotedblright{}.) Say it once more.
\end{tcolorbox}
\section[\texorpdfstring{\ta{வலுப்படுத்து} (valuppaṭuttu)}{வலுப்படுத்து (valuppaṭuttu)}]{\ta{வலுப்படுத்து} (valuppaṭuttu) — to strengthen}

\label{lesson:TA-C390-valuppatuttu}



\begin{quote}
Before the new one: say the Tamil for to improve, then the Tamil for to restructure.
\end{quote}

\subsection*{You'll want to know: \ta{வலுப்படுத்து}}
\textbf{\ta{வலுப்படுத்து}} — \emph{valuppaṭuttu} — \textquotedblleft{}to strengthen\textquotedblright{}.

\begin{quote}
\ta{தினமும்} \ta{உடற்பயிற்சி} \ta{செய்தால்} \ta{அது} \ta{உடலை} \ta{வலுப்படுத்தும்}.
\end{quote}

\begin{quote}
\emph{tiṉamum uṭaṟpayiṟci ceytāl atu uṭalai valuppaṭuttum.}
\end{quote}

\begin{quote}
If you exercise every day, it will strengthen the body.
\end{quote}

Past (I …): \ta{வலுப்படுத்தினேன்}, \emph{valuppaṭuttiṉēṉ}.

What is strengthened takes -\ta{ஐ}. \ta{வலு}, strength.

\begin{cousinweb}
\textbf{\ta{வலுப்படுத்து}} (\emph{valuppaṭuttu}) is native: \emph{valu}, \textquotedblleft{}strength, firmness\textquotedblright{} (DEDR 5276), with \emph{paṭuttu}, the causative of \emph{paṭu}: to make strong.
\end{cousinweb}

\begin{grammarlens}[title={What you've built}]
One more word for saying what happens if.
\end{grammarlens}

\subsection*{Your turn}
\begin{itemize}
\raggedright
  \item \emph{Say it:} \emph{valuppaṭuttu}
  \item \emph{Say it:} \emph{valuppaṭuttu}, once more
  \item \emph{Say it:} \emph{tiṉamum uṭaṟpayiṟci ceytāl atu uṭalai valuppaṭuttum.} — If you exercise every day, it will strengthen the body.
  \item \emph{Recall:} say the Tamil for to improve, then the Tamil for to restructure, then say \emph{valuppaṭuttu} again
\end{itemize}

\begin{tcolorbox}[breakable,colback=teal!4,colframe=teal!35!black,title={Before you move on}]
What does \ta{வலுப்படுத்து} mean? (\textquotedblleft{}To strengthen\textquotedblright{}.) Say it once more.
\end{tcolorbox}
